% ═══════════════════════════════════════════════════════════════
%  chapters/ch11-gtk4-basics.tex
%  Основы GTK4: новый API, виджеты, сигналы
% ═══════════════════════════════════════════════════════════════

\chapter{Основы GTK4}
\label{ch:gtk4-basics}

\epigraph{<<GTK4~--- это не <<ещё один API>>, а \emph{другой}
способ думать о виджетах. Проще, но непривычнее>>.}{}

\section{Введение}
\label{sec:gtk4-intro}

В \texttt{ch01}--\texttt{ch10} мы работали с \emph{железом}:
STM32, ADC, DMA, DAC. \emph{Теперь}~--- \emph{другая} сторона:
\textbf{ПК}. \emph{Программа}, которая \emph{управляет}
осциллографом и генератором.

\textbf{GTK4}~--- это \emph{современная} версия GTK.
В \emph{первой} книге мы использовали \textbf{GTK3}. GTK4~---
\emph{другой}. \emph{Некоторые} вещи \emph{проще},
\emph{некоторые}~--- \emph{сложнее}.

\begin{ideabox}
\textbf{Аналогия: переезд в новую квартиру.} Вещи \emph{те}
же, но \emph{расположение} \emph{другое}. \emph{Поначалу}
\emph{непривычно}, потом~--- \emph{удобнее}.
GTK4~--- \emph{новая} \emph{квартира} для \emph{старых}
\emph{виджетов}.
\end{ideabox}

\subsection{Что вы узнаете}

\begin{itemize}
    \item Чем GTK4 \emph{отличается} от GTK3.
    \item Как \emph{запустить} приложение через \texttt{GtkApplication}.
    \item Как \emph{создавать} виджеты (\texttt{box}, \texttt{button}, \texttt{label}).
    \item Как \emph{подключать} сигналы (по-новому).
    \item Как \emph{рисовать} через \texttt{GtkDrawingArea} и \texttt{Cairo}.
\end{itemize}

\section{GTK4 vs GTK3}
\label{sec:gtk4-vs-gtk3}

\textbf{Главные} \emph{изменения}:

\begin{table}[htbp]
\centering
\caption{GTK3 vs GTK4}
\label{tab:gtk4-vs-gtk3}
\small
\begin{tabularx}{\textwidth}{@{}l L L@{}}
\toprule
\textbf{Что} & \textbf{GTK3} & \textbf{GTK4} \\
\midrule
Запуск       & \texttt{gtk\_init()} & \texttt{GtkApplication} \\
Главный цикл & \texttt{gtk\_main()} & \texttt{g\_application\_run()} \\
Добавление   & \texttt{gtk\_box\_pack\_start()} & \texttt{gtk\_box\_append()} \\
Показ        & \texttt{gtk\_widget\_show\_all()} & \texttt{gtk\_widget\_set\_visible()} \\
Рисование    & \texttt{"draw"} + \texttt{cairo\_t*} & \texttt{gtk\_drawing\_area\_set\_draw\_func()} \\
Стили        & \texttt{GtkStyleContext} & \texttt{GtkCssProvider} (частично) \\
\bottomrule
\end{tabularx}
\end{table}

\begin{keybox}
\textbf{Главное изменение~--- \texttt{GtkApplication}.} Вместо
\emph{глобального} \texttt{gtk\_init()}~--- \emph{объект}
приложения. Это \emph{позволяет} \emph{интегрироваться}
с системой (меню, уведомления, D-Bus).
\end{keybox}

\section{GtkApplication: новый запуск}
\label{sec:gtk4-application}

\textbf{Старый} \emph{способ} (GTK3):
\begin{listing}[htbp]
\begin{minted}[frame=none, bgcolor=gray!5]{c}
int main(int argc, char *argv[])
{
    gtk_init(&argc, &argv);
    GtkWidget *win = gtk_window_new(GTK_WINDOW_TOPLEVEL);
    gtk_widget_show_all(win);
    gtk_main();
    return 0;
}
\end{minted}
\caption{Запуск в GTK3}
\label{lst:gtk3-main}
\end{listing}

\textbf{Новый} \emph{способ} (GTK4):
\begin{listing}[htbp]
\begin{minted}[frame=none, bgcolor=gray!5]{c}
static void on_activate(GtkApplication *app, gpointer data)
{
    GtkWidget *win = gtk_application_window_new(app);
    gtk_window_set_title(GTK_WINDOW(win), "Пример");
    gtk_widget_set_visible(win, TRUE);
}

int main(int argc, char *argv[])
{
    GtkApplication *app = gtk_application_new(
        "org.example.App", G_APPLICATION_DEFAULT_FLAGS);
    g_signal_connect(app, "activate",
                     G_CALLBACK(on_activate), NULL);
    int status = g_application_run(G_APPLICATION(app),
                                   argc, argv);
    g_object_unref(app);
    return status;
}
\end{minted}
\caption{Запуск в GTK4}
\label{lst:gtk4-main}
\end{listing}

\textbf{Что здесь важно}:
\begin{itemize}
    \item \texttt{gtk\_application\_new()}~--- \emph{создаёт}
          \emph{объект} приложения.
    \item \texttt{"activate"}~--- \emph{сигнал}, который
          \emph{вызывается} при \emph{запуске}.
    \item \texttt{gtk\_application\_window\_new(app)}~---
          \emph{создаёт} \emph{окно}, \emph{привязанное}
          к приложению.
    \item \texttt{g\_application\_run()}~--- \emph{главный}
          \emph{цикл}.
\end{itemize}

\begin{ideabox}
\textbf{Зачем \texttt{GtkApplication}?} Чтобы \emph{приложение}
было \emph{уникальным}. Если \emph{запустить} \emph{второй}
\emph{экземпляр}~--- \emph{он} \emph{передаст}
\emph{управление} \emph{первому}, а \emph{не} \emph{откроет}
\emph{новое} \emph{окно}.
\end{ideabox}

\section{Виджеты}
\label{sec:gtk4-widgets}

Виджеты \emph{те} \emph{же}, но \emph{API} \emph{другое}.

\subsection{Box}

\textbf{GTK3}:
\begin{minted}[frame=none, bgcolor=gray!5]{c}
GtkWidget *box = gtk_box_new(GTK_ORIENTATION_VERTICAL, 5);
gtk_box_pack_start(GTK_BOX(box), label, FALSE, FALSE, 0);
gtk_box_pack_start(GTK_BOX(box), button, FALSE, FALSE, 0);
\end{minted}

\textbf{GTK4}:
\begin{minted}[frame=none, bgcolor=gray!5]{c}
GtkWidget *box = gtk_box_new(GTK_ORIENTATION_VERTICAL, 5);
gtk_box_append(GTK_BOX(box), label);
gtk_box_append(GTK_BOX(box), button);
\end{minted}

\begin{keybox}
\textbf{Главное изменение:} вместо \texttt{pack\_start}
\emph{с} \emph{флагами}~--- просто \texttt{append}. \emph{Меньше}
\emph{параметров}, \emph{проще} \emph{код}.
\end{keybox}

\subsection{Button}

\textbf{GTK3}:
\begin{minted}[frame=none, bgcolor=gray!5]{c}
GtkWidget *btn = gtk_button_new_with_label("OK");
g_signal_connect(btn, "clicked",
                 G_CALLBACK(on_clicked), data);
\end{minted}

\textbf{GTK4}:
\begin{minted}[frame=none, bgcolor=gray!5]{c}
GtkWidget *btn = gtk_button_new_with_label("OK");
g_signal_connect(btn, "clicked",
                 G_CALLBACK(on_clicked), data);
\end{minted}

\begin{notebox}
\textbf{Сигналы \emph{не} \emph{изменились}.} \texttt{"clicked"}
\emph{работает} \emph{так} \emph{же}. \emph{Но} \emph{есть}
\emph{новые} \emph{сигналы} (например, \texttt{"activate"}).
\end{notebox}

\subsection{Label}

\textbf{GTK3}:
\begin{minted}[frame=none, bgcolor=gray!5]{c}
GtkWidget *lbl = gtk_label_new("Hello");
gtk_label_set_text(GTK_LABEL(lbl), "World");
\end{minted}

\textbf{GTK4}:
\begin{minted}[frame=none, bgcolor=gray!5]{c}
GtkWidget *lbl = gtk_label_new("Hello");
gtk_label_set_text(GTK_LABEL(lbl), "World");
\end{minted}

\begin{notebox}
\textbf{Label \emph{не} \emph{изменился}.} \emph{Всё} \emph{то}
\emph{же} \emph{API}. \emph{Это}~--- \emph{хорошо}.
\end{notebox}

\section{Сигналы и колбэки}
\label{sec:gtk4-signals}

Сигналы~--- \emph{те} \emph{же}. \emph{Но} \emph{есть}
\emph{новые}:
\begin{itemize}
    \item \texttt{"activate"}~--- \emph{приложение}
          \emph{запущено}.
    \item \texttt{"close-request"}~--- \emph{окно}
          \emph{закрывается}.
    \item \texttt{"destroy"}~--- \emph{окно}
          \emph{уничтожено}.
\end{itemize}

\begin{keybox}
\textbf{Главное} \emph{изменение} \emph{в} \emph{колбэках}~---
\texttt{GtkWidget*} \emph{заменён} \emph{на} \emph{более}
\emph{конкретные} \emph{типы} (\texttt{GtkButton*},
\texttt{GtkWindow*}). \emph{Это}~--- \emph{лучше}
\emph{типизация}.
\end{keybox}

\section{GtkDrawingArea и Cairo}
\label{sec:gtk4-drawing}

\textbf{GTK3}:
\begin{minted}[frame=none, bgcolor=gray!5]{c}
g_signal_connect(area, "draw",
                 G_CALLBACK(on_draw), data);

gboolean on_draw(GtkWidget *w, cairo_t *cr, gpointer data)
{
    cairo_set_source_rgb(cr, 1, 1, 1);
    cairo_paint(cr);
    return FALSE;
}
\end{minted}

\textbf{GTK4}:
\begin{minted}[frame=none, bgcolor=gray!5]{c}
gtk_drawing_area_set_draw_func(
    GTK_DRAWING_AREA(area), on_draw, data, NULL);

void on_draw(GtkDrawingArea *area, cairo_t *cr,
             int width, int height, gpointer data)
{
    cairo_set_source_rgb(cr, 1, 1, 1);
    cairo_paint(cr);
}
\end{minted}

\begin{keybox}
\textbf{Главное} \emph{изменение}: \emph{вместо} \emph{сигнала}
\texttt{"draw"}~--- \emph{функция}
\texttt{gtk\_drawing\_area\_set\_draw\_func()}. \emph{И}
\emph{колбэк} \emph{получает} \emph{размеры} \emph{виджета}
(\texttt{width}, \texttt{height}).
\end{keybox}



%%%%%%========================



\section{Полный пример: окно с кнопкой}
\label{sec:gtk4-example}

\textbf{Теперь}~--- \emph{полный} \emph{пример}. \emph{Окно},
\emph{кнопка}, \emph{label}. \emph{Нажимаем} \emph{кнопку}~---
\emph{меняется} \emph{текст}.

\subsection{Код}

\begin{listing}[htbp]
\begin{minted}[frame=none, bgcolor=gray!5]{c}
#include <gtk/gtk.h>

static void on_button_clicked(GtkButton *btn, gpointer data)
{
    GtkLabel *label = GTK_LABEL(data);
    static int count = 0;
    char buf[64];
    snprintf(buf, sizeof(buf), "Нажато: %d", ++count);
    gtk_label_set_text(label, buf);
}

static void on_activate(GtkApplication *app, gpointer data)
{
    GtkWidget *win = gtk_application_window_new(app);
    gtk_window_set_title(GTK_WINDOW(win), "Пример GTK4");
    gtk_window_set_default_size(GTK_WINDOW(win), 400, 300);

    GtkWidget *box = gtk_box_new(GTK_ORIENTATION_VERTICAL, 10);
    gtk_widget_set_margin_top(box, 20);
    gtk_widget_set_margin_bottom(box, 20);
    gtk_widget_set_margin_start(box, 20);
    gtk_widget_set_margin_end(box, 20);
    gtk_window_set_child(GTK_WINDOW(win), box);

    GtkWidget *label = gtk_label_new("Нажмите кнопку");
    gtk_box_append(GTK_BOX(box), label);

    GtkWidget *button = gtk_button_new_with_label("Нажми меня");
    g_signal_connect(button, "clicked",
                     G_CALLBACK(on_button_clicked), label);
    gtk_box_append(GTK_BOX(box), button);

    gtk_widget_set_visible(win, TRUE);
}

int main(int argc, char *argv[])
{
    GtkApplication *app = gtk_application_new(
        "org.example.Gtk4Demo", G_APPLICATION_DEFAULT_FLAGS);
    g_signal_connect(app, "activate",
                     G_CALLBACK(on_activate), NULL);
    int status = g_application_run(G_APPLICATION(app),
                                   argc, argv);
    g_object_unref(app);
    return status;
}
\end{minted}
\caption{Полный пример GTK4: окно, кнопка, label}
\label{lst:gtk4-example}
\end{listing}

\textbf{Что здесь важно}:
\begin{itemize}
    \item \texttt{gtk\_window\_set\_child()}~--- \emph{новый}
          \emph{способ} \emph{добавить} \emph{ребёнка}
          \emph{в} \emph{окно}.
    \item \texttt{gtk\_box\_append()}~--- \emph{добавить}
          \emph{виджет} \emph{в} \emph{box}.
    \item \texttt{gtk\_widget\_set\_margin\_*}~--- \emph{отступы}
          \emph{от} \emph{краёв}.
    \item \texttt{gtk\_widget\_set\_visible(win, TRUE)}~---
          \emph{показать} \emph{окно}.
\end{itemize}

\subsection{Сборка}

\begin{listing}[htbp]
\begin{minted}[frame=none, bgcolor=gray!5]{bash}
gcc -o demo demo.c `pkg-config --cflags --libs gtk4`
./demo
\end{minted}
\caption{Компиляция и запуск}
\label{lst:gtk4-build}
\end{listing}

\begin{keybox}
\textbf{Команда \texttt{pkg-config}.} \texttt{--cflags}~---
\emph{флаги} \emph{компиляции} (\texttt{-I...}). \texttt{--libs}~---
\emph{флаги} \emph{линковки} (\texttt{-lgtk-4 ...}). \emph{Если}
\emph{не} \emph{хотите} \emph{писать} \emph{вручную}~---
\emph{используйте} \texttt{pkg-config}.
\end{keybox}

\section{Код для нашего проекта}
\label{sec:gtk4-project}

\textbf{Теперь}~--- \emph{код}, \emph{который} \emph{будет}
\emph{в} \emph{нашем} \emph{проекте}. \emph{Это}~---
\emph{заготовка} \emph{для} \emph{осциллографа}.

\subsection{Структура}

\begin{listing}[htbp]
\begin{minted}[frame=none, bgcolor=gray!5]{c}
typedef struct {
    GtkWidget *window;
    GtkWidget *drawing_area;
    GtkWidget *status_label;
    GtkWidget *btn_start;
    GtkWidget *btn_stop;
    // ...
} AppData;
\end{minted}
\caption{Структура \texttt{AppData}}
\label{lst:gtk4-appdata}
\end{listing}

\textbf{Что здесь важно}:
\begin{itemize}
    \item \texttt{AppData}~--- \emph{всё} \emph{состояние}
          \emph{приложения}.
    \item \texttt{drawing\_area}~--- \emph{место} \emph{для}
          \emph{осциллограммы}.
    \item \texttt{status\_label}~--- \emph{статус} \emph{внизу}.
\end{itemize}

\subsection{Активация}

\begin{listing}[htbp]
\begin{minted}[frame=none, bgcolor=gray!5]{c}
static void on_activate(GtkApplication *app, gpointer data)
{
    AppData *st = data;

    st->window = gtk_application_window_new(app);
    gtk_window_set_title(GTK_WINDOW(st->window),
                         "STM32 Oscilloscope");
    gtk_window_set_default_size(GTK_WINDOW(st->window),
                                800, 600);

    GtkWidget *vbox = gtk_box_new(GTK_ORIENTATION_VERTICAL, 5);
    gtk_window_set_child(GTK_WINDOW(st->window), vbox);

    // Кнопки
    GtkWidget *hbox = gtk_box_new(GTK_ORIENTATION_HORIZONTAL, 5);
    gtk_box_append(GTK_BOX(vbox), hbox);

    st->btn_start = gtk_button_new_with_label("Старт");
    g_signal_connect(st->btn_start, "clicked",
                     G_CALLBACK(on_start_clicked), st);
    gtk_box_append(GTK_BOX(hbox), st->btn_start);

    st->btn_stop = gtk_button_new_with_label("Стоп");
    g_signal_connect(st->btn_stop, "clicked",
                     G_CALLBACK(on_stop_clicked), st);
    gtk_box_append(GTK_BOX(hbox), st->btn_stop);

    // Область рисования
    st->drawing_area = gtk_drawing_area_new();
    gtk_widget_set_size_request(st->drawing_area, -1, 300);
    gtk_drawing_area_set_draw_func(
        GTK_DRAWING_AREA(st->drawing_area),
        on_draw, st, NULL);
    gtk_box_append(GTK_BOX(vbox), st->drawing_area);

    // Статус
    st->status_label = gtk_label_new("Готов");
    gtk_box_append(GTK_BOX(vbox), st->status_label);

    gtk_widget_set_visible(st->window, TRUE);
}
\end{minted}
\caption{\texttt{on\_activate()} для проекта}
\label{lst:gtk4-activate}
\end{listing}

\subsection{Рисование}

\begin{listing}[htbp]
\begin{minted}[frame=none, bgcolor=gray!5]{c}
static void on_draw(GtkDrawingArea *area, cairo_t *cr,
                    int width, int height, gpointer data)
{
    AppData *st = data;

    // Фон
    cairo_set_source_rgb(cr, 0.1, 0.1, 0.1);
    cairo_paint(cr);

    // Сетка
    cairo_set_source_rgb(cr, 0.3, 0.3, 0.3);
    cairo_set_line_width(cr, 0.5);
    for (int i = 0; i < width; i += 50) {
        cairo_move_to(cr, i, 0);
        cairo_line_to(cr, i, height);
    }
    for (int i = 0; i < height; i += 50) {
        cairo_move_to(cr, 0, i);
        cairo_line_to(cr, width, i);
    }
    cairo_stroke(cr);

    // Сигнал (пример)
    cairo_set_source_rgb(cr, 0.2, 1.0, 0.2);
    cairo_set_line_width(cr, 2);
    for (int i = 0; i < width; i++) {
        double x = i;
        double y = height / 2.0
                 + 100.0 * sin(i * 0.05);
        if (i == 0) cairo_move_to(cr, x, y);
        else        cairo_line_to(cr, x, y);
    }
    cairo_stroke(cr);
}
\end{minted}
\caption{\texttt{on\_draw()}~--- рисование осциллограммы}
\label{lst:gtk4-draw}
\end{listing}

\textbf{Что здесь важно}:
\begin{itemize}
    \item \texttt{cairo\_set\_source\_rgb()}~--- \emph{цвет}.
    \item \texttt{cairo\_paint()}~--- \emph{заливка} \emph{фона}.
    \item \texttt{cairo\_move\_to()} / \texttt{cairo\_line\_to()}~---
          \emph{линии} \emph{сетки}.
    \item \texttt{cairo\_stroke()}~--- \emph{нарисовать}.
    \item \texttt{sin()}~--- \emph{пример} \emph{сигнала}.
\end{itemize}

\begin{keybox}
\textbf{Почему \texttt{sin(i * 0.05)}?} Это \emph{простой}
\emph{способ} \emph{нарисовать} \emph{синус} \emph{без}
\emph{данных}. \emph{Позже}~--- \emph{заменим} \emph{на}
\emph{реальные} \emph{данные} \emph{из} \emph{порта}.
\end{keybox}

\section{Уроки главы}
\label{sec:gtk4-lessons}

\begin{enumerate}
    \item \textbf{GTK4~--- \emph{другая} \emph{квартира}.}
          Вещи \emph{те} \emph{же}, \emph{API} \emph{другое}.
    \item \textbf{Запуск~--- через \texttt{GtkApplication}.}
          \emph{Уникальность} \emph{приложения}.
    \item \textbf{Виджеты~--- \emph{проще}.} \texttt{append}
          \emph{вместо} \texttt{pack\_start}.
    \item \textbf{Сигналы~--- \emph{те} \emph{же}.}
          \emph{Но} \emph{новые} \emph{типы}.
    \item \textbf{Рисование~--- через
          \texttt{set\_draw\_func}}. \emph{Колбэк}
          \emph{получает} \emph{размеры}.
    \item \textbf{\texttt{AppData}~--- состояние}. \emph{Всё}
          \emph{в} \emph{одной} \emph{структуре}.
    \item \textbf{Cairo~--- тот} \emph{же}. \emph{Методы}
          \emph{не} \emph{изменились}.
\end{enumerate}

\section{Чек-лист грабель}
\label{sec:gtk4-pitfalls}

\begin{table}[htbp]
\centering
\caption{Грабли GTK4}
\label{tab:gtk4-pitfalls}
\small
\begin{tabularx}{\textwidth}{@{}c L L@{}}
\toprule
\textbf{№} & \textbf{Грабля} & \textbf{Решение} \\
\midrule
1 & \texttt{gtk\_init} не работает & \texttt{GtkApplication} \\
2 & Нет окна & \texttt{gtk\_widget\_set\_visible} \\
3 & \texttt{pack\_start} не найден & \texttt{gtk\_box\_append} \\
4 & \texttt{show\_all} не работает & \texttt{set\_visible} \\
5 & Нет рисования & \texttt{set\_draw\_func} \\
6 & Нет размера в draw & \emph{Колбэк} \emph{получает} \texttt{width}, \texttt{height} \\
7 & Нет \texttt{pkg-config} & \texttt{--cflags --libs gtk4} \\
8 & Не видно окна & \emph{Проверить} \texttt{set\_visible} \\
\bottomrule
\end{tabularx}
\end{table}

\section{Упражнения}
\label{sec:gtk4-exercises}

\begin{exercisebox}
\textbf{Упражнение 1.} Соберите \emph{пример} из
\cref{lst:gtk4-example}. \emph{Запустите}. \emph{Нажмите}
\emph{кнопку}. \emph{Что} \emph{произошло}?
\end{exercisebox}

\begin{exercisebox}
\textbf{Упражнение 2.} Добавьте \emph{вторую} \emph{кнопку}
(<<Сброс>>). \emph{При} \emph{нажатии}~--- \emph{счётчик}
\emph{обнуляется}. \emph{Как} \emph{связать} \emph{кнопки}?
\end{exercisebox}

\needspace{5\baselineskip}
\begin{exercisebox}
\textbf{Упражнение 3.} Измените \emph{цвет} \emph{фона}
\emph{и} \emph{сетки} \emph{в} \texttt{on\_draw()}.
\emph{Как} \emph{сделать} \emph{сетку} \emph{более}
\emph{заметной}?
\end{exercisebox}

\begin{exercisebox}
\textbf{Упражнение 4.} Добавьте \emph{слайдер}
(\texttt{GtkScale}) \emph{для} \emph{частоты}. \emph{При}
\emph{изменении}~--- \emph{меняется} \emph{частота}
\emph{сигнала}. \emph{Какой} \emph{сигнал} \emph{использовать}?
\end{exercisebox}

\section{Проверьте сами}
\label{sec:gtk4-check}

\begin{enumerate}
    \item Соберите \emph{пример} \emph{из}
          \cref{lst:gtk4-example}.
    \item \emph{Запустите}. \emph{Нажмите} \emph{кнопку}.
    \item \emph{Убедитесь}, что \emph{счётчик} \emph{работает}.
    \item \emph{Измените} \emph{цвет} \emph{фона} \emph{на} \emph{белый}.
    \item \emph{Сделайте} \emph{скриншот}~--- \emph{пришлите}
          \emph{в} \emph{книгу}.
\end{enumerate}

\begin{ideabox}
\textbf{Если получилось~--- вы поняли GTK4.} Это
\emph{главное}. Дальше~--- \emph{потоки} \emph{и} \emph{данные}
\emph{из} \emph{порта}. \emph{Самое} \emph{интересное}~---
\emph{впереди}.
\end{ideabox}


